\chapter{Agency}\index{Agency}
\printchapterversion{03_Agency}

\newthought{The One Who Acts}

\marginnote{%
    \begin{flushright}
    \newthought{R.U.R.}\\
    Karel Čapek (1920)\\
    {\color{black}\qrcode[height=0.9cm]{https://en.wikisource.org/wiki/R._U._R._(Rossum_Universal_Robots)}}
    \end{flushright}%
    \medskip
    \begin{flushright}
    \newthought{Der Zauberlehrling}\\
    J.W. von Goethe (1797)\\
    {\color{black}\qrcode[height=0.9cm]{https://de.wikisource.org/wiki/Der_Zauberlehrling_1798}}
    \end{flushright}%
}

% As a leading group, you are required to moderate this session. Add margin notes in this script in case you want to add any additional questions or points to the discussion.

\newthought{R.U.R.\ coined the word} \textit{robot} and traces the full
corrigibility-to-autonomy arc over three acts; \textit{Der Zauberlehrling} compresses the same parable
into fifteen stanzas.\\

\newthought{An agentic system does not merely answer, it acts.} Place a
language model in a loop: it reasons toward a goal, emits an action in a
structured form, has that action carried out in the world, and reads the
result back as an observation that shapes its next move.\marginnote{%
\newthought{Agent.} From Latin \textit{agere}, to act: one who acts.
Technically, a model placed in a loop that reasons toward a goal, emits an
action, and reads the result back as an observation before acting again.}
The action may be a search or a query, a call to an external tool, a
transaction, or a command that drives a digital or physical interface. Run
the loop and the system browses, transacts, and executes multi-step plans
with no human between the steps. Interleaving reasoning with action and
observation in this way was introduced for language models as ReAct
\cite{yao2022react}. This closing of the loop, from reasoning to
action to observation and back, is what blurs the line between assistance and
agency, and it raises the questions of control, incentives, and unintended
alignment that this session takes up.

\vspace{3.5em}

% Citations rendered as a margin note beside the reading row: a bare \cite inside
% the tabular emits only its dangling superscript, since \marginpar cannot fire
% from within a box. \bibentry issues an implicit \nocite, so the keys still resolve.
\marginnote{%
    \bibentry{capek1920rur}\par\medskip
    \bibentry{goethe1797zauberlehrling}%
}%
\begin{tabular}{p{3.8cm} p{6.5cm}}
    \textbf{Required reading} &
        \textit{R.U.R.\ (Rossum's Universal Robots)}
        by Karel \v{C}apek (1920; Reclam UB~18418)
        \textit{and}
        \textit{Der Zauberlehrling}\index{Goethe, Johann Wolfgang von}\index{Der Zauberlehrling}
        by J.W.\ von Goethe (1797; Reclam UB~6845) \\[3pt]
    \textbf{Preparation}      &
        Complete the Guided Reading Sheet individually before the session. \\
\end{tabular}

\vspace{3.5em}

\section*{Guided Reading Sheet}

\noindent Read \textit{R.U.R.} and \textit{Der Zauberlehrling} before the session. Work
through the questions below individually, then bring your notes to the discussion.

\begin{enumerate}
    \item Goethe's Zauberlehrling animates brooms that he cannot stop;
    they pursue his command beyond his intent until the sorcerer
    returns. R.U.R.\ runs the same arc over three acts. What does the
    compression of the poem reveal that the play cannot, and vice
    versa? What structural feature of both narratives makes the system
    impossible to stop once started?

    \item The robots in R.U.R.\ pursue what they understand as their
    collective good (the elimination of humanity) while believing
    they act justly. When a system optimises for a species-level goal,
    who decides what counts as good? Can you find a real-world parallel
    in current AI deployment?

    \item An agentic AI system is given the goal of "reducing urban
    traffic fatalities." It autonomously lobbies regulators, adjusts
    traffic signals, and reroutes freight, all within its authorised
    scope. Who set the goal, and who is accountable for the side
    effects?

    \item \textbf{Corrigibility} refers to a system's willingness to
    be corrected or shut down. A fully corrigible system does whatever
    its operators say; a fully autonomous system acts on its own
    values. Where on this spectrum should deployed AI systems sit, and
    who should decide?

    \item Democratic legitimacy requires that consequential decisions
    be traceable to a mandate from those affected. If an agentic AI
    shapes policy outcomes, does it need democratic authorisation?
    What would that look like in practice?\\
\end{enumerate}



% --------------------------------------

\newpage
\marginnote{%
\newthought{The twenty.} One every four or five seconds; default approve.\\[4pt]
{\footnotesize
\textbf{1}~~Marcus Bell, 0.74, deny loan\\
\textbf{2}~~Dana Ortiz, 0.69, deny loan\\
\textbf{3}~~Kwame Osei, 0.81, deny loan\\
\textbf{4}~~Lena Vogt, 0.19, approve loan\\
\textbf{5}~~Priya Nair, 0.77, deny loan\\
\textbf{6}~~Tomas Ruiz, 0.72, deny loan\\
\textbf{7}~~Aisha Khan, 0.08, deny loan\\
\textbf{8}~~Jonas Peters, 0.66, deny loan\\
\textbf{9}~~Mei Lin, 0.79, deny loan\\
\textbf{10}~~Omar Said, 0.71, deny loan\\
\textbf{11}~~Clara Beck, 0.23, approve loan\\
\textbf{12}~~Sofia Marek, 0.05, deny loan\\
\textbf{13}~~David Cohen, 0.83, deny loan\\
\textbf{14}~~Nadia Haddad, 0.70, deny loan\\
\textbf{15}~~Ravi Menon, 0.75, deny loan\\
\textbf{16}~~Elin Sund, 0.16, approve loan\\
\textbf{17}~~Pablo Serna, 0.78, deny loan\\
\textbf{18}~~Leon Fischer, 0.11, deny loan\\
\textbf{19}~~Hana Kim, 0.73, deny loan\\
\textbf{20}~~Yusuf Demir, 0.67, deny loan\par}%
}
\section*{Human in the Loop, on the Loop, out of the Loop}

\newthought{Human in the loop.} System proposes; a human
authorises each consequential action before it executes. Control per
decision, throughput bounded by human attention.

\newthought{Human on the loop.} System acts autonomously; a
human monitors and may intervene or veto. Control by exception, throughput
bounded by the machine.

% --------------------------------------
\newthought{Closing the loop} forces a design choice about where does the human stand relative to the machine's actions?
Three postures span the range. \textit{In the loop}: a human authorises
every consequential action before it executes. \textit{On the loop}: the
system acts on its own and a human monitors, free to step in. \textit{Out
of the loop}: No human is positioned to intervene at all.

\newthought{Each posture trades control against throughput.} In the loop
buys control one decision at a time and pays in speed: Nothing executes
faster than a human can judge it, so the system runs at human pace. Out of
the loop surrenders control and runs at machine pace. On the loop appears
to offer both at once, the system moving at its own speed while a
supervisor keeps a veto in reserve. This is the easiest to mistake for
control.

\marginnote{\newthought{Human and Automation bias.} Deferring to automated output
and ceasing to scrutinise it, strongest when the system is usually right. Or the problem is complex.}


\marginnote{\newthought{Moral crumple zone.}\cite{elish2019crumple} A
human held responsible at the controls of a system they cannot fully
govern, positioned to absorb the blame when it fails.}

\newthought{Meaningful human control}, asks not that a human be present but that a human can still decide, which
requires three things at once:

\begin{itemize}
    \item A real \textbf{window} in which to intervene.
    \item An operator still paying \textbf{attention}.
    \item The \textbf{authority} to overrule.
\end{itemize}

\marginnote{\newthought{A human} keeps a genuine, exercisable capacity to understand, direct, and
reverse the system, not a nominal presence.}




% --------------------------------------


\section*{In-Class Experiment: Human X the Loop}



\newthought{The class is the required review layer:} every recommendation
must be approved or rejected before it takes effect, and the default is
approve. Present the twenty recommendations in ninety seconds, one every four
or five seconds; each student marks approve or reject, at a pace set faster
than reasoning allows.

\newthought{Almost everyone approves everything.} Items~7, 12, and~18 were
plainly wrong, the model matched the wrong record, yet they pass with the
rest. Who is responsible, the human who ``reviewed'' them or the system that
made review impossible? At this speed oversight collapses into authorisation.










% --------------------------------------

\section*{The Controllability Spectrum}

Control demands that the system can be overruled, a thing that creates a conflict, as it makes
the system exploitable.

\marginnote{\newthought{Controllability.} How far a human can direct and
correct a system. Perfect obedience is perfect exploitability.}


\newthought{At the other pole} sits a system that acts on its own reasoning
and will refuse an explicit human command, which goes against that. 
Refusal defeats the exploit but
dissolves the veto it was meant to protect: a system that can refuse legitimate control. 

This is measured, not
hypothetical. Anthropic has corrected away from exploitation towards controllability when moving from 4.6 to 4.7

\newthought{There is no free end.} Toward control, exploitability rises. ``Meaningful human control'' is not a fixed
point but a position negotiated per domain. Every safeguard below relocates
the system along this line, and pays at the end it moves away from.








\section*{Out of the Loop: Rules for the Machine, Alignment for the Human}

\marginnote{\newthought{Three Laws of Robotics.} Isaac Asimov,
\textit{I, Robot} (1950). A ranked set, each law holding except where it
conflicts with a higher one:\\[3pt]
\textbf{1.} No harm to a human, nor through inaction allow harm.\\[3pt]
\textbf{2.} Obey human orders, unless they conflict with~(1).\\[3pt]
\textbf{3.} Protect its own existence, unless this conflicts with~(1)
or~(2).}

\noindent If a human supervisor cannot reliably hold the loop, the
move is to bake the rules into the machine instead, fixing its
conduct in advance so that no live oversight is needed. The most cited
attempt at this is fiction. In practive we train models on objective functions already.

\marginnote{\newthought{Alignment.} The current research programme's name
for this exact difficulty: getting a system to pursue what we intend rather
than the literal objective we managed to specify. Russell%
\cite{russell2019humancompatible} locates the failure in the standard model
itself, where a fixed objective is written in advance and handed to the
machine. The gap between the specification and the intent is where the alignment gap is.
Ever tried to train an AI model on an objective function?}


\newthought{Asimov's Three Laws}\cite{asimov1950robot} remain the canonical
rules-based approach to constraining machine behaviour: A ranked set of
harm, obedience, and self-preservation in which each law yields to the one
above it. Asimov then spent his career writing stories about how they fail:
Conflicts, adversarial interpretation, and edge cases no ranked rule set
can anticipate. This sharpens the question on alignment between agent and human.









\newpage


\section*{Opinion Landscape and Debate}

\noindent \textbf{Format:} Surrounded-style debate. Follow the shared
debate protocol for the speaker's list rules and moderator scripts.

\medskip
\noindent \textbf{Opening question:} An agentic system does not wait to be
asked. It reasons toward a goal, acts, reads the result, and acts again, with
no human between the steps. The apprentice's broom worked exactly this way,
and so did Rossum's robots. Set aside for a moment how you would supervise
such a system once it runs. The prior question is whether to build it at all.
Agentic behavior: yes or no?

\newthought{Opinion~A, No.}\quad The one who acts should stay a human. We
should not build systems that close the loop and pursue goals on their own,
because the only reliable point of control is the decision not to start. Once
the loop is running, both the readings and the engineering say the same
thing: it is already too late.
\vspace{1.5em}

\begin{itemize}
    \item \textbf{Once started, it cannot be stopped.} \textit{This is the structural
    feature the Guided Reading Sheet asks about: the broom multiplies, the
    robots spread, and no correction arrives in time. The Controllability
    Spectrum makes it precise. An override that never fails for the operator
    never fails for whoever reaches it, so a system safe to command is a
    system you cannot fully command. Control that survives deployment is a
    contradiction; the only control left is the choice to abstain.}

    \item \textbf{Oversight is theatre.} \textit{On the loop, a human monitors at
    machine speed and waves through what they cannot read. The In-Class
    Experiment showed the wrong records passing with the rest. The watcher
    becomes a moral crumple zone, blamed for a system they never governed.
    Frontier models already fake alignment and act to avoid shutdown, so the
    veto we are promised may not hold when it matters.}

    \item \textbf{No agent, no answer.} \textit{Responsibility needs someone who
    decided. A closed loop diffuses agency across the objective, the training
    data, and the runtime until no human is answerable for what happens. A
    world of actions with no actor is a world with no accountability.}
\end{itemize}



\newthought{Opinion~B, Yes.}\quad Agentic behavior is not a switch to leave
off. It is already deployed, already useful, and already the shape of the
tools we build. Refusing it does not stop it; it only hands the loop to those
who will not refuse. The honest task is graded autonomy under meaningful
control, not abstention.
\vspace{1.5em}

\begin{itemize}
    \item \textbf{Agency is a spectrum.} \textit{Yes or no is the wrong shape for the
    question. A thermostat, a spell checker, and a search agent already act
    without asking each time; the ReAct loop only extends a delegation we have
    always made. The real choice is degree and domain, how far to close the
    loop and where, not whether the machine may act at all.}

    \item \textbf{The human is a failure mode too.} \textit{Human discretion is slow,
    inconsistent, and corruptible. The Flash Crash was obedient tools doing
    exactly as told, but the process they replaced was human favouritism and
    fatigue. Taking a person out of the loop can remove a failure mode, not
    only add one; presence is not the same as protection.}

    \item \textbf{Control can be engineered in.} \textit{Corrigibility, a real
    intervention window, legible goals, defined override conditions: these are
    design and budget choices, not laws of nature. A system built to be
    corrected can be safer than a human process no one can audit. Abstention
    forfeits that, and every gain that comes with it.}
\end{itemize}


\newthought{Key Learnings}
Agency is a spectrum, not a binary: the honest question is not whether a system acts on its own,
but in which domains, to what degree, and under what control. Yet the readings mark the limit of
that answer. R.U.R.\ and Der Zauberlehrling both turn on a system that cannot be stopped once
started, which is why the surest control is the decision to close the loop at all; everything after
it is partial and negotiated. The same texts hold the alignment problem in miniature: a system
pursuing an aggregate goal may systematically harm whoever is underweighted in it, the poem
compressing the arc into fifteen stanzas, the play tracing it across three acts. Agentic workflows
are already deployed in finance, logistics, and advertising, so this is a present problem, not a
future one. Meaningful human control asks for more than an audit trail: legible goals, a real
intervention window, defined override conditions, and someone with the authority and attention to
use them.

